\documentclass[11pt,a4paper]{article}

\usepackage[T1]{fontenc}
\usepackage[utf8]{inputenc}
\usepackage[margin=22mm]{geometry}
\usepackage{booktabs}
\usepackage{array}
\usepackage{longtable}
\usepackage{enumitem}
\usepackage{titlesec}
\usepackage{fancyhdr}
\usepackage{listings}
\usepackage[hidelinks]{hyperref}

\pagestyle{fancy}
\fancyhf{}
\lhead{\small BeamMP Race Manager}
\rhead{\small Phase 3 result and phase 4 plan}
\cfoot{\small \thepage}
\renewcommand{\headrulewidth}{0.4pt}

\titleformat{\section}{\large\bfseries}{\thesection}{0.6em}{}
\titlespacing{\section}{0pt}{14pt}{6pt}

\setlist[itemize]{leftmargin=*,itemsep=2pt,topsep=4pt}
\setlist[enumerate]{leftmargin=*,itemsep=2pt,topsep=4pt}
\setlength{\parskip}{6pt}
\setlength{\parindent}{0pt}

\lstset{
  basicstyle=\ttfamily\scriptsize,
  frame=single,
  framesep=5pt,
  columns=fullflexible,
  keepspaces=true,
  breaklines=true,
  showstringspaces=false,
  xleftmargin=0pt,
  aboveskip=8pt,
  belowskip=8pt,
}

\newcommand{\code}[1]{\texttt{\small #1}}

\begin{document}

\begin{center}
{\LARGE\bfseries Phase 3 is done, and what phase 4 is}\\[4pt]
{\small BeamMP Race Manager \quad Client: dhardesty99}
\end{center}

\vspace{4pt}
\hrule
\vspace{8pt}

Phase 3 is finished, including the three changes you asked for after testing.
One thing you asked for cannot be done, and section 4 says why in full.

Phase 4 was going to be racing together on one grid. Your two new ideas, the
pit checkpoint and starting a race properly, belong in it, so they are in here
as part of the same phase.

\section{What the bottom bar does now}

\begin{longtable}{@{}p{24mm}p{12mm}p{16mm}p{72mm}@{}}
\toprule
\textbf{Button} & \textbf{Hold} & \textbf{Penalty} & \textbf{What it does} \\
\midrule
\endhead
Reposition & 5s & 60s  & Puts the car back on its wheels where it stands. \\
Spare tire & 30s & 30s & Fixes the car where it stands. Only when a tire is flat. \\
Repair     & 60s & 30s & Fixes the car where it stands. \\
Fuel +25\% & 20s & none & Adds a quarter of a tank. \\
Lights     & none & none & Head lights, fog, hazards, lightbar, siren, horn, flash. \\
\bottomrule
\end{longtable}

\textbf{Hold} is how long the car is held still. \textbf{Penalty} is seconds
added to your race time. The clock keeps running during a hold, so a repair in
a race costs a minute of waiting and thirty seconds on the time.

Outside a run there is no hold and no penalty. The buttons just work. This is
worth knowing when you test: on the grid, before you cross the line, you are
still outside the run, so everything is instant and free.

Every number lives on the server and can change without anyone reinstalling
the mod.

\section{The three things you found}

All three were the same mistake. Each one moved the car.

\begin{longtable}{@{}p{26mm}p{58mm}p{58mm}@{}}
\toprule
& \textbf{What was wrong} & \textbf{What it does now} \\
\midrule
\endhead
Reposition
& It teleported you to the last checkpoint, and it also repaired the car for
  free, so there was never a reason to press Repair.
& Sets the car back on its wheels where it stands. Nothing is repaired. No
  checkpoint is needed. \\
\addlinespace
Repair
& It put the car at the last place it had stopped.
& Repairs where it stands. \\
\addlinespace
Spare tire
& The same teleport as Repair.
& No teleport. See section 4 for the part that cannot be fixed. \\
\bottomrule
\end{longtable}

Reposition was repairing because of one setting. The game call that moves a car
takes an argument that decides whether to mend it on the way, and it is on
unless you turn it off. Repair was moving the car because it was calling the
recovery function, which is the one that puts you back where you were last
stopped. It uses the reset call now, which fixes the car and leaves it where it
is.

\section{The spare tire, and the one thing the game will not do}

You asked for the spare tire to change only the broken tire, in place. In place
is done. Only the broken tire cannot be.

BeamNG can burst a tire and it cannot mend one. Its own code offers
\code{deflateTire}, \code{deflateTires} and \code{deflateRandomTire}, and
nothing that goes the other way. The only place the game clears the flat flag is
inside a full reset of every wheel. There is no call that repairs one wheel and
leaves the rest of the car alone, so there is nothing for me to reach for.

So Spare tire fixes the whole car, like Repair does. What keeps the two buttons
apart is that Spare tire is only allowed when a tire is actually flat. You
cannot use it as a cheap repair, because you need a flat first.

Three ways to go. My default is the first.

\begin{enumerate}
  \item Leave it. Spare tire is the cheaper repair you can only use with a flat.
  \item Make it cost the same as Repair, so the choice is only about the wait.
  \item Drop Spare tire and keep Repair.
\end{enumerate}

The tire shop sound is a separate thing and it is easy. It needs a sound file
from you, or I can use one of the game's own.

\section{Phase 4: racing together}

Today everybody drives the same course alone and the results are collected when
the last car is off track. Phase 4 makes it one race.

\begin{itemize}
  \item One grid. Everyone is placed in a row behind the start line.
  \item One set of results, in finishing order, with gaps to the leader and to
        the car in front.
  \item Qualifying, then a race, if you want the two split up.
  \item Somebody leaving no longer holds up everyone else.
\end{itemize}

\section{Phase 4: public and invited races}

This is your idea and it replaces the part you called touchy.

Right now a run arms the moment you pick a course, and the clock starts the
first time you touch the start line. That is easy to trigger by accident.

Instead, a race becomes a thing you make, and other people join.

\begin{enumerate}
  \item You pick a course, the laps and the mode, and you make a race.
  \item You choose \textbf{public}, so anyone on the server can join, or
        \textbf{invite}, and you pick who.
  \item People join. Everyone who joins is put on the grid.
  \item Nothing is armed. The line does nothing. You can sit there as long as
        you like.
  \item You press \textbf{Start Race}. Only then does the line become live.
  \item Each driver's clock starts when they cross the first checkpoint.
\end{enumerate}

So the start line still starts your own time, which keeps the timing honest,
but it does nothing at all until the host says go.

\begin{lstlisting}
local lobbies = {}

function RM.lobby.create(pid, d)
  local track = RM.tracks.get(tostring(d.track or ""))
  if not track then return false, "no_such_track" end
  if RM.lobby.of(pid) then return false, "already_in_one" end

  local id = RM.util.shortId()
  lobbies[id] = {
    id      = id,
    host    = pid,
    track   = track.id,
    laps    = math.floor(tonumber(d.laps) or 1),
    mode    = tostring(d.mode or "controller"),
    open    = d.open and true or false,
    state   = "waiting",
    members = { [pid] = true },
    invited = {},
  }
  return true, lobbies[id]
end

function RM.lobby.start(pid)
  local l = RM.lobby.of(pid)
  if not l then return false, "no_lobby" end
  if l.host ~= pid then return false, "not_the_host" end
  if l.state ~= "waiting" then return false, "already_started" end

  l.state = "armed"
  for member in pairs(l.members) do
    RM.race.arm(member, { id = l.track, mode = l.mode, laps = l.laps })
    RM.bus.queue(member, "race.state", RM.race.wire(member))
  end
  return true, l
end
\end{lstlisting}

\section{Phase 4: the pit checkpoint}

Your other idea. A checkpoint can be marked as a pit when the course is built.

A pit is not part of the lap. Driving through it does not count as a gate and
does not skip one. What it does is put you in the pit while you are inside it.

While you are in the pit:

\begin{itemize}
  \item Fuel fills to 100\% instead of adding a quarter.
  \item Repair and Spare tire add nothing to your race time.
  \item The holds still apply in full. A repair still costs you the minute of
        waiting, and the clock is still running.
\end{itemize}

So a pit stop costs you the time it takes and nothing on top. That is the point
of having one.

\begin{lstlisting}
function RM.service.use(pid, d)
  local action = ACTIONS[tostring(d.which or "")]
  if not action then return false, "no_such_action" end
  if jobs[pid] then return false, "already_working" end

  local r = RM.race.get(pid)
  local racing = r ~= nil and r.state == "running"
  local inPit = racing and r.inPit == true

  local hold, cost = 0, nil
  if racing then
    hold = tonumber(RM.config.holds[action.hold]) or 0
    if action.penalty and not inPit then
      local seconds = tonumber(RM.config.penalties[action.penalty])
      if seconds and seconds > 0 then
        RM.race.penalty(pid, seconds, action.penalty)
        cost = seconds
      end
    end
  end

  jobs[pid] = { which = action.name, hold = hold, endsAt = RM.now() + hold,
                cost = cost, full = inPit }
  return true, jobs[pid]
end
\end{lstlisting}

The course builder gets one more control: mark the checkpoint you just dropped
as a pit. Pit checkpoints are drawn in a different colour so they are obvious.

\section{Order of work}

\begin{enumerate}
  \item The pit checkpoint. It is small, it is self contained, and it is
        useful on its own.
  \item Lobbies on the server: make, join, invite, start. All testable without
        the game.
  \item The race panel: a list of open races, a join button, an invite list and
        the Start Race button.
  \item One grid for everybody, and one set of results.
  \item Qualifying, if you want it split from the race.
\end{enumerate}

The server side comes first every time, because most of it can be proven before
the game is even opened. That is what worked in phases 2 and 3.

\section{Questions, and what I will do if you say nothing}

\begin{longtable}{@{}p{6mm}>{\raggedright\arraybackslash}p{70mm}>{\raggedright\arraybackslash}p{72mm}@{}}
\toprule
& \textbf{Question} & \textbf{Default if you say nothing} \\
\midrule
\endhead
1 & The spare tire cannot fix only one tire. Which of the three in section 4?
  & The first. It stays the cheaper repair that needs a flat first. \\
\addlinespace
2 & Can people join a race after it has started?
  & No. Once you press Start Race the field is closed. \\
\addlinespace
3 & What happens to a driver who joins and then never moves?
  & The race carries on. They are filed as a did not finish, the same as now. \\
\addlinespace
4 & Can there be more than one race going at once?
  & Yes, on different courses. One race per course at a time. \\
\addlinespace
5 & Who can make a race public?
  & Anybody. Invites are picked from the player list. \\
\addlinespace
6 & Does a pit checkpoint have to be driven through slowly?
  & No speed rule for now. There is a 37 mph limit sitting unused in the
    settings if you want one later. \\
\addlinespace
7 & Should the pit also fix a flat for free?
  & Yes. In the pit nothing adds to your time. \\
\addlinespace
8 & The tire shop sound. Your file or one of the game's?
  & One of the game's, until you send one. \\
\bottomrule
\end{longtable}

\section{Not in phase 4}

Records and personal bests are phase 5. CoPilot and Challenges are phase 6.

The look of the interface is a separate job and can happen at any point, as we
said. It does not need to wait for phase 6.

\vspace{10pt}
\hrule
\vspace{6pt}

{\small Phase 3 is on the server now. Phase 4 starts on your word, and the eight
questions above are the only thing that would change the shape of it.}

\end{document}
